\sgaStableTarget{sga6.e12.expose}\section*{\underline{EXPOSÉ XII}\\[3pt] \underline{\underline{A RELATIVE REPRESENTABILITY THEOREM FOR}}\\[3pt] \underline{\underline{THE PICARD FUNCTOR}}}
\begin{center}
\underline{by M. Raynaud (written up by S. Kleiman)}
\end{center}

\sgaStableTarget{sga6.e12.warning}\subsection*{Warning}

The purpose of the two exposés XII, XIII is to give proofs of the nine finiteness theorems for the Picard functor stated by A. Grothendieck in the comments (pp. C 07-C 11) to his ``Fondements de la Géométrie Algébrique'' (cited as FGA). In the first exposé, the first finiteness theorem is proved in the following form, more precise than the one stated in FGA: let @@M0@@ and @@M1@@ be two proper schemes over an integral noetherian base @@M2@@, and let @@M3@@ be a surjective @@M4@@-morphism. Then there is a nonempty open subset @@M5@@ of @@M6@@ over which the induced morphism @@M7@@ is representable by a quasi-affine morphism of finite presentation. In the second exposé, the rest is proved.

The two exposés differ notably in their proof techniques. In the first, essential use is made of Oort's dévissage, non-flat descent, and the representation of unramified functors. In the second, one first reduces to the case where the schemes @@M8@@ that occur are projective and flat with geometrically integral fibers, or even, as needed, normal or smooth, always by applying the first finiteness theorem to a suitable morphism @@M9@@. One then uses the tool of @@M10@@-sheaves forged by G. Castelnuovo and sharpened and used by T. Matsusaka, D. Mumford, Y. Nakai and S. Kleiman. This tool makes it possible, in particular, to dispense with Chow coordinates, the Frobenius morphism and the Lefschetz theorems used in Grothendieck's initial proofs.

The present exposé XII was presented orally, following Grothendieck, by M. Raynaud, whom the writer wishes to thank for communicating his personal notes.

